\chapter{Stato dell'arte nel rilevamento automatico delle cadute}
\label{cap:stato_arte}

% ---------------------------------------------------------------------
% RISCRITTO IL 2026-08-25. Terza stesura.
%
% Prima stesura: 5.271 parole, venticinque rimandi ai capitoli
% sperimentali, diciannove paragrafi con citazione su centosette.
% Seconda: 2.645 parole, ma otto sottosezioni, sei delle quali dentro
% la stessa sezione e lunghe fra 104 e 189 parole. Un titolo ogni
% duecento parole.
%
% Questa stesura applica due criteri.
%
% RILEVANZA. Sono state tolte le parti che descrivono cose che la tesi
% non impiega: il confronto fra elaborazione diretta dell'immagine e
% posa (la scelta e' argomentata nel Cap. 1 sulla riservatezza, non
% qui), e la trattazione estesa delle reti a grafo, mai provate in
% questo lavoro. Di entrambe resta una menzione, perche' un lettore che
% conosce il settore si aspetta di trovarle, ma non una sezione.
%
% GERARCHIA. Le sei sottosezioni della sezione sui dati sono diventate
% \paragraph: hanno un nome ma non una voce nell'indice, che e' il
% trattamento giusto per un blocco da centocinquanta parole.
% ---------------------------------------------------------------------

Il Capitolo~\ref{cap:contesto} ha definito il problema applicativo e i
requisiti del sistema. Questo capitolo richiama le soluzioni proposte in
letteratura e i limiti dei protocolli con cui vengono valutate, allo scopo di
fissare i concetti e la terminologia impiegati nel resto della trattazione.

L'ordine seguito è quello delle scelte da compiere: il sensore, che determina
quale informazione sia disponibile; la rappresentazione, che stabilisce quale
parte di essa raggiunga il classificatore; il modello temporale; infine i dati
e i protocolli di valutazione.

\section{Tecnologie per il rilevamento delle cadute}
\label{sec:famiglie}

Una classificazione consolidata distingue i sistemi di rilevamento in tre
famiglie: dispositivi indossabili, sensori ambientali e sistemi basati su
visione artificiale~\cite{mubashir, gutierrez}.

I sistemi \emph{indossabili} applicano il sensore al corpo. Accelerometri e
giroscopi sono le tecnologie più diffuse, perché misurano direttamente
accelerazione, orientamento e velocità angolare; le prime soluzioni
impiegavano regole a soglia, quelle successive classificatori supervisionati su
finestre temporali dei segnali inerziali. Il vantaggio è che il sensore segue
la persona e non dipende dal punto di vista di una telecamera. Il limite è la
dipendenza dalla \emph{compliance} dell'utente: il dispositivo va indossato
correttamente, mantenuto carico e usato con continuità~\cite{gazem}.

I sensori \emph{ambientali}, quali pressione, acustica, infrarossi e radar, non
richiedono nulla all'utente, ma trasferiscono il vincolo all'ambiente: il
sistema copre soltanto le aree strumentate e richiede installazione o
calibrazione. Molte di queste tecnologie osservano inoltre il movimento in modo
indiretto, e forniscono una descrizione più povera di un flusso visivo.

I sistemi \emph{vision-based} impiegano telecamere RGB, RGB-D o termiche. Una
rassegna dedicata a questa famiglia ha analizzato ottantuno sistemi pubblicati
fra il 2015 e il 2020, documentando la transizione da descrittori geometrici
progettati a mano verso rappresentazioni apprese~\cite{gutierrez}. La visione
fornisce più informazione di molti sensori ambientali: da una sequenza si
stimano posizione, configurazione corporea, dinamica e interazione con
l'ambiente. Comporta però due costi: l'elaborazione di interi fotogrammi
richiede memoria e capacità di calcolo, e il video contiene informazioni
identificative del soggetto e dell'abitazione.

Il presente lavoro adotta la visione, e affronta il secondo costo con la scelta
architetturale descritta nella sezione~\ref{sec:privacy}.

\section{Rappresentazione del movimento}
\label{sec:rappresentazioni-sota}

Il sensore stabilisce quale informazione esista; la rappresentazione stabilisce
quale parte di essa raggiunga il classificatore.

I primi metodi vision-based impiegavano caratteristiche geometriche della
sagoma segmentata: rapporto fra altezza e larghezza del rettangolo di
contenimento, inclinazione, posizione del baricentro, variazione dell'area
occupata. Con la diffusione dell'apprendimento profondo l'immagine è diventata
essa stessa l'ingresso del modello, e la rappresentazione viene appresa dai
dati. In entrambi i casi la descrizione resta però legata all'apparenza della
scena, e quindi a sfondo, illuminazione e disposizione dell'ambiente.

La stima della posa trasforma invece un'immagine in un insieme di punti
anatomici: il classificatore riceve una descrizione della configurazione
corporea e non del contenuto visivo. È la scelta adottata in questo lavoro,
per le ragioni di riservatezza esposte nella sezione~\ref{sec:privacy}.

Gli approcci basati sullo scheletro sono ampiamente studiati nel
riconoscimento delle azioni umane. Una rassegna del 2023, che raccoglie circa
duecentocinquanta risultati da oltre settanta studi pubblicati fra il 2019 e il
marzo 2023, distingue le caratteristiche geometriche progettate a mano dalle
rappresentazioni apprese~\cite{nguyen}. La distinzione individua due strategie
alternative: i punti possono essere impiegati direttamente come coordinate,
lasciando che sia il modello a scoprire quali relazioni contino, oppure
trasformati in grandezze progettate che descrivano esplicitamente angolo del
tronco, altezze relative e velocità. Il Capitolo~\ref{cap:metodologia}
costruisce entrambe.

Alla rappresentazione scheletrica vengono attribuite due proprietà utili a un
sistema destinato a hardware limitato. La prima è la \emph{compattezza}: poche
decine di valori sostituiscono centinaia di migliaia di componenti RGB. La
seconda è l'\emph{interpretabilità}: ogni variazione del segnale si riconduce a
una parte del corpo, il che rende possibile misurare separatamente il
contributo dei singoli gruppi di punti.

La posa non è però una misura ma una stima, e risente di occlusioni,
illuminazione, risoluzione e scorci prospettici; le coordinate bidimensionali
dipendono inoltre dal punto di vista. Una catena basata sulla posa sposta
quindi parte della difficoltà dal classificatore allo stimatore e alla
geometria costruita sopra di esso.

Nella fall detection questa impostazione è già stata impiegata con modelli
sequenziali: Lin et al. estraggono la posa con OpenPose e confrontano LSTM e
GRU per modellarne l'evoluzione temporale~\cite{lin}.

\section{Modellazione temporale e classificatori}
\label{sec:modelli-sota}

Fissata la rappresentazione, il rilevamento diventa un problema di
classificazione di sequenze, e la letteratura propone quattro famiglie
principali: ricorrenti, convoluzionali, a grafo e basate su
attenzione~\cite{nguyen}.

Le reti \emph{ricorrenti} sono la scelta naturale quando l'ingresso è una
sequenza. LSTM e GRU introducono meccanismi di memoria e di controllo del
flusso pensati per conservare informazione nel tempo; GRU adotta una struttura
più compatta e, a parità di dimensione dello stato, richiede meno parametri.
Le \emph{convoluzioni temporali} apprendono schemi locali all'interno della
sequenza e non richiedono uno stato aggiornato campione per campione. Le reti
\emph{a grafo} sfruttano la struttura naturale dello scheletro, in cui i punti
sono i nodi e le connessioni anatomiche gli archi. I modelli \emph{ad
attenzione} assegnano peso differente ai vari istanti e mettono in relazione
diretta elementi anche distanti nella sequenza.

A queste va aggiunto un termine di paragone non sequenziale. Nei sistemi
basati su grandezze progettate la conoscenza del fenomeno è incorporata nella
rappresentazione e non nel modello, e classificatori quali le foreste casuali
hanno costo ridotto e permettono di esaminare l'importanza delle variabili. Un
classificatore che non modella l'ordine dei campioni è inoltre necessario per
stabilire se un modello sequenziale stia effettivamente sfruttando dipendenze
temporali.

Questo lavoro confronta reti ricorrenti, convoluzioni temporali, un modello ad
attenzione e un classificatore non sequenziale; le reti a grafo non vengono
considerate.

La progressione verso architetture più espressive rende plausibile attendersi
un effetto rilevante della scelta del modello. I confronti pubblicati sono
però difficili da interpretare, perché variano contemporaneamente dati,
grandezze, pre-elaborazione, divisione degli insiemi e soglia decisionale.
Isolare l'effetto della sola architettura richiede un protocollo in cui il
resto della catena resti invariato.

\section{Dati e protocolli di valutazione}
\label{sec:dati-sota}

La raccolta di cadute reali presenta vincoli etici e pratici evidenti: non è
possibile provocare deliberatamente incidenti in persone fragili per costruire
un archivio. Gran parte dei riferimenti pubblici contiene quindi cadute
simulate da volontari giovani e sani.

La simulazione non è pertanto una scorciatoia metodologica, ma l'unica via
disponibile per ottenere il numero di esempi positivi che un classificatore
richiede. La domanda pertinente è un'altra: che cosa, di preciso, la
simulazione non conserva. Per la componente dinamica la risposta è
documentata. Kangas et al. hanno confrontato cadute accidentali reali di
persone anziane, raccolte con sensori indossati per sei mesi, con cadute
sperimentali eseguite da soggetti di mezza età: la grandezza più indicativa
risulta essere la rapidità di crescita della decelerazione anziché il suo
picco, e nelle cadute reali compaiono impatti multipli, assenti negli altri
movimenti~\cite{kangasreali}.

\paragraph{Gli archivi pubblici.} Fra quelli impiegati nella letteratura
vision-based compaiono UR-Fall, Le2i, Multiple Cameras Fall Dataset e UP-Fall.
Differiscono per numero di soggetti, attività, tipi di caduta, punti di vista,
frequenza di ripresa e criteri di annotazione, e l'eterogeneità rende privo di
significato il confronto diretto fra valori di $F$-score ottenuti su archivi
diversi. UP-Fall, adottato in questo lavoro, comprende diciassette soggetti
giovani e sani, undici attività, più prove per attività e due telecamere
sincronizzate~\cite{upfall}: la presenza di due viste consente di studiare
l'effetto della geometria di ripresa, e la varietà delle attività quotidiane
permette di costruire negativi più realistici di quelli disponibili in archivi
composti quasi soltanto da cadute e posture erette.

\paragraph{Verificabilità.} Fra gli ottantuno sistemi recensiti da Gutiérrez
et al. la maggioranza è stata valutata su dati propri, per i quali non
risultava accesso pubblico al momento della revisione~\cite{gutierrez}. Un
risultato ottenuto su cadute simulate resta confrontabile se l'archivio è
disponibile; un risultato ottenuto su un archivio privato non è né
confrontabile né riproducibile.

\paragraph{Prevalenza e negativi difficili.} Negli archivi le cadute devono
essere abbastanza frequenti da permettere addestramento e valutazione, mentre
nella vita quotidiana la classe positiva è rara; la
sezione~\ref{sec:falsiallarmi} ne ha svolto il conto. Il rapporto fra negativi
e positivi non è tuttavia il criterio con cui giudicare la composizione di un
archivio, perché nessun archivio potrebbe riprodurre il rapporto reale: farlo
richiederebbe di registrare anni di vita ordinaria per ogni caduta osservata.
Il criterio pertinente riguarda \emph{quali} negativi siano presenti. Diventano
allora centrali i \emph{negativi difficili}, cioè le attività ordinarie che
condividono caratteristiche con la classe positiva: sedersi di scatto,
inginocchiarsi, chinarsi e distendersi volontariamente producono configurazioni
parzialmente simili a quelle di una caduta. Se tali attività sono assenti o
sottorappresentate, il modello può apprendere scorciatoie che funzionano
sull'archivio e generano falsi allarmi in esercizio.

\paragraph{Azioni e stati.} Un passo recente verso una valutazione uniforme è
OmniFall, che unifica otto archivi allestiti in una tassonomia comune a sedici
classi, quindicimila riprese per ottanta ore con annotazione per
fotogramma~\cite{omnifall}. Il vocabolario distingue le azioni
\emph{transitorie}, come \texttt{fall} e \texttt{lie\_down}, dagli \emph{stati}
statici, come \texttt{fallen} e \texttt{lying}. Su un insieme di verifica
costituito da video di incidenti reali gli autori riscontrano che i modelli
valutati senza adattamento continuano a confondere lo stato \texttt{fallen} con
\texttt{lying}, mentre quelli addestrati su scene contenenti esplicitamente la
permanenza a terra ottengono risultati sensibilmente migliori. La confusione
documentata non riguarda quindi la distinzione fra caduta e attività
quotidiana, ma quella fra una persona a terra e una persona distesa.

\paragraph{Tre livelli di generalizzazione.} Anche all'interno dello stesso
archivio il modo in cui si costruiscono gli insiemi determina il significato
del risultato. Finestre sovrapposte estratte dalla stessa ripresa sono
fortemente correlate: distribuite a caso fra addestramento e verifica, portano
a valutare il modello su sequenze quasi identiche a quelle già osservate. La
divisione \emph{per soggetto} impone invece di generalizzare verso persone
assenti dall'addestramento, e protocolli come il \emph{leave-one-subject-out}
rendono visibile la dispersione fra individui che una media aggregata
nasconde. Per i sistemi vision-based esiste un secondo livello, la
generalizzazione fra \emph{punti di vista}, poiché una grandezza
bidimensionale può essere informativa da una vista laterale e degenerare
quando lo stesso movimento è osservato dall'alto. Il terzo livello, il
trasferimento fra \emph{archivi}, è il più severo, perché cambiano insieme
persone, ambiente, telecamera, frequenza di ripresa, annotazione e
distribuzione delle attività.

\paragraph{Dalla finestra all'evento.} La maggior parte dei classificatori
temporali produce una decisione per ogni finestra, mentre un sistema installato
deve produrre \emph{eventi} di allarme, e le due quantità non coincidono: una
situazione che genera dieci finestre positive consecutive può essere aggregata
in un unico evento, e una finestra isolata può essere soppressa da una regola
di persistenza. La rassegna di Broadley et al. sui metodi di valutazione in
condizioni reali rileva che le prove sul campo sono intrinsecamente difficili,
che gli studi disponibili impiegano insiemi piccoli, e che manca accordo su
come quantificare la capacità di evitare i falsi allarmi e su che cosa
costituisca un non-evento~\cite{broadley}. Precisione, richiamo e $F$-score
calcolati sulla finestra restano perciò necessari ma non sufficienti.

\paragraph{Esecuzione su dispositivo.} Molti lavori riportano accuratezza e
complessità del classificatore senza misurare l'intera catena su hardware
reale, benché il costo comprenda acquisizione, decodifica, stima della posa,
costruzione delle grandezze, classificazione e logica temporale. Una mappatura
sistematica del 2026 raccoglie quattrocentotrentatré studi primari pubblicati
fino al 2025 e vi dedica una parte specifica: fotogrammi al secondo, latenza e
piattaforma hardware effettivamente impiegata~\cite{gomes}. La velocità di
esecuzione non è quindi una proprietà dell'architettura, ma una misura da
compiere sul dispositivo di destinazione e sulla catena completa.

\section{Questioni aperte e posizionamento del lavoro}
\label{sec:posizionamento}

Dalla letteratura emerge un quadro nel quale esistono molte soluzioni
tecnicamente valide, e restano aperte quattro questioni che diventano rilevanti
nel passaggio da un archivio di riferimento a un sistema destinato
all'utilizzo reale.

La prima riguarda la \emph{rappresentazione del movimento}: immagine, sagoma,
coordinate scheletriche e grandezze progettate presentano vantaggi e limiti
diversi quanto a interpretabilità, costo computazionale e dipendenza dal punto
di vista~\cite{gutierrez, gomes}.

La seconda riguarda il \emph{modello di classificazione}. Le famiglie
richiamate nella sezione~\ref{sec:modelli-sota} sono tutte strategie plausibili
per lo stesso problema~\cite{nguyen, lin}, ma i risultati riportati in lavori
differenti sono difficili da confrontare, perché variano contemporaneamente
dati, rappresentazione e protocollo~\cite{omnifall}.

\begin{Ques}
Quanto la scelta dell'architettura influenza il risultato, quando
rappresentazione, dati e protocollo vengono mantenuti costanti?
\end{Ques}

La terza riguarda il \emph{significato della classe positiva}: riconoscere la
transizione di caduta, riconoscere lo stato successivo a terra e decidere
quando generare un allarme sono problemi collegati ma non
equivalenti~\cite{omnifall}.

La quarta riguarda la \emph{generalizzazione}, poiché cambiamenti di soggetto,
punto di vista, ambiente e archivio modificano il comportamento della catena,
così come il passaggio dall'esecuzione fuori linea a un dispositivo
reale~\cite{broadley, gomes}.

\begin{Ques}
Che cosa un archivio di laboratorio non riesce a misurare, e con quali
conseguenze per chi lo impiega per scegliere fra alternative?
\end{Ques}

\bigskip

Il presente lavoro si colloca all'interno di queste questioni. Viene adottata
una rappresentazione basata sulla posa, sulla quale sono costruite e
confrontate diverse descrizioni geometriche e temporali del movimento; più
famiglie di classificatori vengono valutate mantenendo invariati, per quanto
possibile, dati e protocollo. La valutazione non si limita al comportamento su
un singolo archivio, ma considera la generalizzazione verso soggetti assenti
dall'addestramento, la variazione del punto di vista, il trasferimento verso un
archivio esterno e il comportamento della catena in condizioni più vicine a
quelle d'uso. Il sistema viene infine considerato nella sua interezza,
includendo il passaggio dalle singole decisioni agli eventi di allarme e la
verifica dell'esecuzione sul dispositivo impiegato.

Il Capitolo~\ref{cap:metodologia} descrive il modo in cui il video viene
trasformato nella rappresentazione utilizzata dal sistema; il
Capitolo~\ref{cap:protocollo} definisce il protocollo con cui le diverse
scelte vengono confrontate.
